\section{Proofs and Auxiliary Results}
\label{app:theory}
\label{app:learning_proof}
\setcounter{equation}{0}

\subsection{Maximum-Sharpe direction}
\label{app:maximum_sharpe}

\begin{proof}[Proof of Proposition~\ref{prop:population_optimal_policy}]
Write $W=aU$, with $\|U\|_{\Hh}=1$ and $a\in\mathbb R$.
For a direction with positive payoff variance, put
$m_U=\E[R^p_{t+1}(U)]$ and $s_U^2=\E[(R^p_{t+1}(U))^2]$.
Then
\[
Q(aU)=1-2am_U+a^2s_U^2,
\qquad
a^\star(U)=\frac{m_U}{s_U^2},
\qquad
\min_a Q(aU)=\frac{1}{1+\SR^2(R^p_{t+1}(U))}.
\]
A zero payoff has loss one; a nonzero constant payoff can be scaled to
have loss zero. The assumption $0<Q(W^\star)<1$ therefore excludes both
from the optimum. Minimizing over directions maximizes squared Sharpe,
and the optimally scaled payoff has mean $m_U^2/s_U^2>0$.
Thus the selected orientation attains the positive maximum Sharpe ratio.
\end{proof}

\subsection{Pricing-kernel interpretation and Hansen--Jagannathan equality}
\label{app:sdf_hj}

\begin{proof}[Proof of Proposition~\ref{prop:sdf_hj}]
Suppress the economy subscript. The first-order condition for
$Q(W^\star+aV)$ at $a=0$ gives
$\E[m^\star R^p(V)]=0$ for every $V\in\Hh$.
Taking $V=W^\star$ and using Proposition~\ref{prop:population_optimal_policy},
\[
\E[m^\star]=\E[(m^\star)^2]=Q(W^\star)
=:q^\star=\frac{1}{1+(\SR^\star)^2}.
\]
Hence $\widetilde m=m^\star/q^\star$ has mean one and
variance $(1-q^\star)/q^\star=(\SR^\star)^2$.
For any square-integrable $d$ with $\E[d]=1$ and
$\E[dR^p(V)]=0$ for every $V\in\Hh$,
$\E[dm^\star]=1$ and therefore
\[
\Var(d)=\Var(\widetilde m)+\E[(d-\widetilde m)^2]
\geq(\SR^\star)^2.
\]
This proves the minimum-variance property without imposing positivity
or the stronger conditional restriction in Condition~\ref{ec:pricing}.
\end{proof}

\subsection{Managed-portfolio Markowitz representation}
\label{app:managed_markowitz}

\begin{proof}[Proof of Proposition~\ref{prop:managed_markowitz}]
The reproducing property gives
$R^p_{t+1}(W)=\langle W,X_t\rangle_{\Hh_K}$ and
$\|X_t\|_{\Hh_K}^2=N^{-2}R_{t+1}'K_tR_{t+1}$.
Assumption~\ref{ass:kernel} gives strong measurability, while
Condition~\ref{ec:returns} and Jensen's inequality give
$\E\|X_t\|_{\Hh_K}^2\leq B_X^2\log2$.
Thus the first and second moments are well defined, and
\[
Q(W)=1-2\langle W,\mu_{\Hh_K}\rangle_{\Hh_K}
+\langle W,\Sigma_{\Hh_K}W\rangle_{\Hh_K}.
\]
The population first-order condition is
$\Sigma_{\Hh_K}W_K^\star=\mu_{\Hh_K}$; the minimum-norm solution is the
stated pseudoinverse representation. The empirical penalized criterion
gives
\[
(\widehat\Sigma_T+\lambda I)\widehat W_{\lambda,T}=\widehat\mu_T.
\]
Since $\widehat\Sigma_T\succeq0$ and $\lambda>0$, the inverse exists and
yields \eqref{eq:markowitz_representation}.
\end{proof}

\subsection{A sufficient Gaussian return specification}
\label{app:primitive_consequences}

\begin{proposition}[A sufficient Gaussian return specification]
\label{prop:gaussian_admissible}
Suppose that, conditionally on $\mathscr F_t$, $R_{t+1}$ is Gaussian, possibly
with a singular covariance, and
\[
\max_{1\leq i\leq N}\E[R_{i,t+1}^2\mid\mathscr F_t]\leq V_R<\infty
\quad\text{a.s.},
\]
uniformly in the economy and time. Then Condition~\ref{ec:returns} holds
with $B_X^2=8\kappa_KV_R$ and $K_4=3$.
\end{proposition}
\begin{proof}
For every $V\in\Hh_K$, $R^p_{t+1}(V)$ is conditionally Gaussian. A scalar
$G\sim N(m,v)$ satisfies
\[
\E G^4=m^4+6m^2v+3v^2\leq3(m^2+v)^2,
\]
which proves \eqref{eq:conditional_fourth_condition} with $K_4=3$.

For the tail condition, Cauchy--Schwarz in $\Hh_K$ and
$K(z,z)\leq\kappa_K$ give
\[
\|X_t\|_{\Hh_K}^2
\leq
\frac{\kappa_K}{N}\sum_{i=1}^N R_{i,t+1}^2.
\]
If $G\sim N(m,v)$ and $m^2+v\leq V_R$, then
\[
\E e^{G^2/(8V_R)}
=\left(1-\frac{v}{4V_R}\right)^{-1/2}
  \exp\!\left(\frac{m^2}{8V_R-2v}\right)
\leq\frac{2}{\sqrt3}e^{1/6}<2.
\]
Therefore, by convexity,
\[
\begin{aligned}
\E_\varepsilon\!\left[
\left.
\exp\!\left(\frac{\|X_t\|^2}{8\kappa_KV_R}\right)
\right|\mathscr F_t
\right]
&\leq
\E_\varepsilon\!\left[
\left.
\exp\!\left(\frac1N\sum_{i=1}^N
              \frac{R_{i,t+1}^2}{8V_R}\right)
\right|\mathscr F_t
\right]\\
&\leq
\frac1N\sum_{i=1}^N
\E_\varepsilon\!\left[
\left.
\exp\!\left(\frac{R_{i,t+1}^2}{8V_R}\right)
\right|\mathscr F_t
\right]
<2.
\end{aligned}
\]
Independence across stocks is unnecessary.
\end{proof}

This is a sufficient specification for Condition~\ref{ec:returns};
Gaussianity is not imposed on the class $\mathcal E$.

\subsection{Exact reduction to the kernel learning problem}
\label{app:induced_krr}

\begin{lemma}[Induced KRR problem]
\label{lem:induced_krr}
For every $\varepsilon\in\mathcal E$, the direct estimator is isometrically
equivalent to scalar KRR with constant response one. The induced problem
satisfies the conditions of \citet[Lemma~A.7]{bignozzigagliardinischneider2026}
with source order $r=1$. Its covariance operator is unitarily equivalent to
$\Sigma_{\Hh_K,\varepsilon}$, and its prediction error is exactly population
portfolio regret.
\end{lemma}

\begin{proof}
\textit{Step 1: identify the learning problem.}
Let $\mathcal X=\mathscr Z^N\times\mathbb R^N$ and
$\Xi_t=(Z_t,R_{t+1})$. For $x=((z_1,\ldots,z_N),r)$, put
\[
\Phi(x)=\frac1N\sum_i r_iK(\cdot,z_i),
\qquad
f_W(x)=\langle W,\Phi(x)\rangle_{\Hh_K}.
\]
The map $\mathcal T:W\mapsto f_W$ is injective: choosing
$r=(N,0,\ldots,0)$ and varying $z_1$ shows that $f_W=0$ on
$\mathcal X$ implies $W=0$. Equip
$\mathcal G=\{f_W:W\in\Hh_K\}$ with
$\langle f_W,f_V\rangle_{\mathcal G}=\langle W,V\rangle_{\Hh_K}$.
Then $\mathcal T$ is an isometry onto the separable RKHS with kernel
\[
\widetilde K(x,x')
=\langle\Phi(x),\Phi(x')\rangle_{\Hh_K}
=\frac1{N^2}\sum_{i,j}r_i s_jK(z_i,z_j'),
\qquad x'=((z_1',\ldots,z_N'),s).
\]
Since $\Phi(\Xi_t)=X_t$, $Y_t\equiv1$ gives exactly the empirical
criterion in \eqref{eq:estimator}, with the same penalty.
Let $J:\mathcal G\to L^2(\mathbb P_{\varepsilon,\Xi})$ be the canonical
embedding. Square integrability, established above, gives
\[
\mathcal T^*J^*J\mathcal T=\Sigma_{\Hh_K,\varepsilon},
\qquad
f_{\mathcal G}=\mathcal TW^\star_\varepsilon.
\]
Expansion of the population criterion at its minimizer gives
\[
\|Jf_W-Jf_{\mathcal G}\|_{L^2}^2
=\|\Sigma_{\Hh_K,\varepsilon}^{1/2}(W-W^\star_\varepsilon)\|_{\Hh_K}^2
=Q_\varepsilon(W)-Q_\varepsilon(W^\star_\varepsilon).
\]

\medskip
\noindent\textit{Step 2: verify the statistical assumptions.}
The induced kernel is measurable, and
$\widetilde K(\Xi_t,\Xi_t)=\|X_t\|_{\Hh_K}^2$.
For the statistical history
$\mathcal A_t=\sigma(\Xi_s,Y_s:s\leq t)$, economic timing gives
$\mathcal A_{t-1}\subseteq\mathscr F_t$.
Condition~\ref{ec:returns} and the tower property therefore imply
\[
\E_\varepsilon\!\left[
\left.\exp\!\left(\frac{\widetilde K(\Xi_t,\Xi_t)}{B_X^2}\right)
\right|\mathcal A_{t-1}\right]\leq2.
\]
This is the conditional kernel-moment assumption in
\citet{bignozzigagliardinischneider2026}, with $B_K=B_X$ and $C_K=2$.
Stationarity and mixing follow from Condition~\ref{ec:dependence};
Condition~\ref{ec:target} gives source order $r=1$, and
Condition~\ref{ec:spectrum} gives the same positive spectrum under the isometry.

It remains to verify the score condition of the cited lemma.
Write $m^\star=m^\star_{t+1,\varepsilon}$,
$P_t^\star=R^p_{t+1}(W^\star_\varepsilon)$ and
$q_t=\E_\varepsilon[(P_t^\star)^2\mid\mathscr F_t]$.
Condition~\ref{ec:pricing} with $V=W^\star_\varepsilon$ gives
$\E_\varepsilon[P_t^\star\mid\mathscr F_t]=q_t$; conditional Jensen
then gives $q_t^2\leq q_t$, so $0\leq q_t\leq1$.
Consequently, Condition~\ref{ec:returns} implies
$\E_\varepsilon[(m^\star)^4\mid\mathscr F_t]\leq8(1+K_4)$.
For every $V\in\Hh_K$, conditional Cauchy--Schwarz yields
\[
\begin{aligned}
\E_\varepsilon[(m^\star)^2(R^p_{t+1}(V))^2\mid\mathscr F_t]
&\leq
\sqrt{\E_\varepsilon[(m^\star)^4\mid\mathscr F_t]
\E_\varepsilon[(R^p_{t+1}(V))^4\mid\mathscr F_t]}\\
&\leq
C_{\rm sc}\E_\varepsilon[(R^p_{t+1}(V))^2\mid\mathscr F_t],
\end{aligned}
\]
where $C_{\rm sc}=\sqrt{8K_4(1+K_4)}$.
Integration and summation over an orthonormal basis show that
$m^\star X_t$ is square integrable and that its second-moment operator
is bounded above by $C_{\rm sc}\Sigma_{\Hh_K,\varepsilon}$.
Condition~\ref{ec:pricing}, separability and the tower property give
$\E_\varepsilon[m^\star X_t\mid\mathcal A_{t-1}]=0$.
The statistical score is $\mathcal T(m^\star X_t)$, so the cited lemma
applies with its variance constant equal to $C_{\rm sc}$.
All constants are common to $\mathcal E$.
\end{proof}

\subsection{Proof of the portfolio-learning rate}
\label{app:rate}

\begin{proof}[Proof of Theorem~\ref{thm:rate}]
Lemma~\ref{lem:induced_krr} verifies the conditions of
\citet[Lemma~A.7]{bignozzigagliardinischneider2026} with $r=1$.
Since $\gamma>b$ and $\lambda_T^\star\asymp T^{-b/(b+1)}$, its upper
rate and the loss identity give \eqref{eq:portfolio-loss-rate} uniformly
over $\mathcal E$.

For the Sharpe statement, put
$D_T=Q_\varepsilon(\widehat W_{\lambda_T^\star,T})-Q_\varepsilon(W^\star_\varepsilon)$,
$q^\star=Q_\varepsilon(W^\star_\varepsilon)$ and
$S^\star=\SR^\star_{\varepsilon,K}$.
Proposition~\ref{prop:population_optimal_policy} and
Condition~\ref{ec:sharpe} give
\[
\underline q:=\frac1{1+\overline S^2}
\leq q^\star\leq\frac1{1+\underline S^2}=: \overline q<1.
\]
On $D_T<1-\overline q$, the estimated loss is below one, hence the
estimated payoff has positive mean. Its variance cannot be zero:
a positive constant payoff could be rescaled to one, contradicting
$q^\star>0$. Write its positive Sharpe as $\widehat S_T$.
Optimizing only its scale gives
\[
q^\star\leq\frac1{1+\widehat S_T^2}\leq q^\star+D_T,
\qquad
0\leq(S^\star)^2-\widehat S_T^2
\leq\frac{D_T}{q^\star(q^\star+D_T)}.
\]
Therefore
\[
0\leq S^\star-\widehat S_T
\leq\frac{D_T}{\underline S\,\underline q^2}.
\]
The event has probability tending to one uniformly, which proves
\eqref{eq:sharpe-level-rate}. Any finite convention on a zero-variance
event leaves this rate unchanged. At a fixed volatility $\sigma$, the
capital-allocation-line gap is $\sigma(S^\star-\widehat S_T)$ and has
the same rate, uniformly for $\sigma$ in a bounded interval.
\end{proof}

\subsection{Effective portfolio complexity}
\label{app:complexity}

\begin{proof}[Proof of Theorem~\ref{thm:effective_complexity}]
Condition~\ref{ec:spectrum} and
\citet[Proposition~A.6]{bignozzigagliardinischneider2026} give
$\mathcal C_{K,\varepsilon}(\lambda)\asymp\lambda^{-1/b}$ uniformly.
At the regularization scale in Theorem~\ref{thm:rate},
\[
\mathcal C^\star_{K,T}\asymp T^{1/(b+1)},
\qquad
\frac{\mathcal C^\star_{K,T}}{T}\asymp T^{-b/(b+1)}\longrightarrow0.
\]
Finiteness follows from
$\mathcal C_{K,\varepsilon}(\lambda)\leq\Tr(\Sigma_{\Hh_K,\varepsilon})/\lambda$.
This is the rate-optimal balance of the loss upper bound, not a shape
theorem for the exact Sharpe curve of every economy.
\end{proof}


\subsection{A finance-specific minimax construction}
\label{app:minimax}

\begin{proof}[Proof of Theorem~\ref{thm:minimax_words}]
The construction keeps the policy kernel fixed. It is a least-favourable family
used only in this lower-bound proof, not a structural return equation imposed on
the economies in the main text.

\medskip
\noindent\textit{Step 1: construct a least-favourable family with nondegenerate Sharpe.}
Since $\mathcal E$ is nonempty, take the one-period state-return marginal
$\nu_0$ of one economy in the class. Write a draw as
$x=(\mathbf z,r)$ and put
\[
Q_X(x):=\|\Phi(x)\|_{\Hh_K}^2,
\qquad
q_X:=\int Q_X(x)\,d\nu_0(x).
\]
Condition~\ref{ec:returns} and Jensen's inequality imply
$q_X\leq B_X^2\log2$. Moreover $q_X>0$, because otherwise the managed
feature would vanish almost surely, contradicting the spectral lower bound
in Condition~\ref{ec:spectrum}.

Define a tilted law by
\[
d\widetilde\nu(x)
=
\frac{Q_X(x)}{q_X}\,d\nu_0(x),
\]
and, on the support of the tilted law, replace the return vector by
\[
r^\circ
=
\sqrt{\frac{q_X}{Q_X(x)}}\,r.
\]
Let $\nu$ denote the law of $(\mathbf z,r^\circ)$ and reuse $r$ for the
normalized return vector below. Since the feature map is linear in returns,
\[
\int \Phi(\mathbf z,r^\circ)\otimes\Phi(\mathbf z,r^\circ)\,d\widetilde\nu
=
\int \Phi(x)\otimes\Phi(x)\,d\nu_0(x)
=:\Sigma_0.
\]
Thus the normalization preserves the entire managed-payoff second-moment
operator. Under $\nu$,
\[
\|\Phi(x)\|_{\Hh_K}^2=q_X\leq B_X^2\log2
\quad\text{almost surely}.
\]
Put $\overline B_X:=B_X\sqrt{\log2}$.

Write
\[
\Sigma_0e_j=\mu_je_j,
\qquad
c_\mu j^{-b}\leq\mu_j\leq C_\mu j^{-b}.
\]
Its eigenvectors are orthonormal in $\Hh_K$. Fix constants $a,h>0$,
with $h\leq a$, sufficiently small that
\[
a^2+h^2\leq R^2,
\qquad
\overline B_X\sqrt{a^2+h^2}\leq\frac12,
\qquad
8C_\mu h^2\leq\frac18.
\]
For an integer $M\geq1$ and $\theta\in\{-1,1\}^M$, define
\[
W_\theta
:=
a e_1
+
\frac{h}{\sqrt M}
\sum_{j=M+1}^{2M}\theta_{j-M}e_j,
\qquad
g_\theta(x)
:=
\langle W_\theta,\Phi(x)\rangle_{\Hh_K}.
\]
Then
\[
\|W_\theta\|_{\Hh_K}=\sqrt{a^2+h^2}\leq R,
\qquad
|g_\theta(x)|\leq \overline B_X\sqrt{a^2+h^2}\leq\frac12.
\]
Put
\[
A_M
:=
\|g_\theta\|_{L^2(\nu)}^2
=
a^2\mu_1
+
\frac{h^2}{M}
\sum_{j=M+1}^{2M}\mu_j.
\]
This quantity does not depend on $\theta$. Moreover,
\[
a^2c_\mu\leq A_M\leq\frac14.
\]

\medskip
\noindent\textit{Step 2: construct admissible financial observations.}
Draw $x_t=(Z_t,r_t)$ independently with law $\nu$. Conditional on $x_t$,
draw $s_t\in\{-1,1\}$ according to
\[
\mathbb P_\theta(s_t=1\mid x_t)
=
\frac{1+g_\theta(x_t)}{2}.
\]
The observed financial row is
$(Z_t,R_{t+1})=(Z_t,s_tr_t)$. For the least-favourable construction take
\[
\mathscr F_t
=
\sigma(x_u,s_u:u<t)\vee\sigma(x_t).
\]
This information structure is allowed by Definition~\ref{def:economy}.
Rows are i.i.d. The managed feature is $s_t\Phi(x_t)$ and, since
$s_t^2=1$,
\[
\|s_t\Phi(x_t)\|_{\Hh_K}^2=q_X.
\]
Because $x_t$ is contained in $\mathscr F_t$,
\[
\E_\theta\!\left[
\left.
\exp\!\left(
\frac{\|s_t\Phi(x_t)\|_{\Hh_K}^2}{B_X^2}
\right)
\right|\mathscr F_t
\right]
=
\exp(q_X/B_X^2)
\leq2.
\]
Moreover, for every $V\in\Hh_K$,
\[
(R^p_{t+1}(V))^2=\langle V,\Phi(x_t)\rangle_{\Hh_K}^2,
\qquad
(R^p_{t+1}(V))^4=\langle V,\Phi(x_t)\rangle_{\Hh_K}^4
\]
conditionally on $\mathscr F_t$. Hence
Condition~\ref{ec:returns} holds with fourth-moment coefficient one,
and therefore with the class constant $K_4\geq1$.
The second-moment operator of the managed feature is exactly $\Sigma_0$.

For every $V\in\Hh_K$, conditional on the current $x_t$ and the past,
\[
\begin{aligned}
&\E_\theta\!\left[
(1-R^p_{t+1}(W_\theta))R^p_{t+1}(V)
\mid x_t,\mathscr F_{t-1}
\right]\\
&\qquad=
\langle V,\Phi(x_t)\rangle_{\Hh_K}
\E_\theta[s_t-g_\theta(x_t)\mid x_t,\mathscr F_{t-1}]
=0.
\end{aligned}
\]
Thus Condition~\ref{ec:pricing} holds. For any portfolio policy $W$ with
finite loss, write $g_W(x)=N^{-1}\sum_iW(z_i)r_i$. Then
\begin{equation}
Q_\theta(W)-Q_\theta(W_\theta)
=
\int(g_W(x)-g_\theta(x))^2\,d\nu(x).
\label{eq:lower_risk_identity}
\end{equation}
Indeed, conditioning on $x$ gives the quadratic
$1-2g_\theta(x)g_W(x)+g_W(x)^2$. Consequently, $W_\theta$ is the
population-optimal minimum-norm policy.

Its payoff satisfies
\[
\E_\theta[R^p_{t+1}(W_\theta)]=A_M,
\qquad
\E_\theta[(R^p_{t+1}(W_\theta))^2]=A_M.
\]
Hence
\[
(\SR_\theta^\star)^2
=
\frac{A_M}{1-A_M}.
\]
Because $A_M$ is uniformly bounded away from zero and from one, the optimal
Sharpe ratio is uniformly bounded away from zero and infinity. Thus the
least-favourable family satisfies Condition~\ref{ec:sharpe}, with constants
common to the family, as well as Conditions~\ref{ec:dependence}--\ref{ec:spectrum}.

\medskip
\noindent\textit{Step 3: neighboring economies are difficult to distinguish.}
Consider the more informative experiment that also reveals $x_t$ and $s_t$.
A lower bound in this augmented experiment also applies to the financial
observations. Neighboring parameters differ in one sign, say coordinate
$j-M$. Their conditional Bernoulli probabilities lie in $[1/4,3/4]$, and
\[
\operatorname{KL}(\operatorname{Ber}(p),\operatorname{Ber}(q))
\leq
\frac{(p-q)^2}{q(1-q)}.
\]
Hence, for the laws $\mathbb P^{\mathrm{aug},T}_\theta$ of $T$ augmented
observations and a one-bit neighbor $\theta'$, we have
\[
\begin{aligned}
\operatorname{KL}(\mathbb P^{\mathrm{aug},T}_\theta,
                  \mathbb P^{\mathrm{aug},T}_{\theta'})
&\leq
2T\int(g_\theta-g_{\theta'})^2\,d\nu\\
&=
\frac{8Th^2\mu_j}{M}\\
&\leq
8C_\mu h^2\frac{T}{M^{b+1}}.
\end{aligned}
\]
Choose
\[
M=\left\lceil T^{1/(b+1)}\right\rceil.
\]
The right-hand side is at most $1/8$, so Pinsker's inequality gives total
variation at most $1/4$ for every neighboring pair.

\medskip
\noindent\textit{Step 4: testing difficulty implies scale-free directional error.}
Define
\[
\psi_j(x)
:=
\frac{\langle e_j,\Phi(x)\rangle_{\Hh_K}}{\sqrt{\mu_j}},
\]
so that $(\psi_j)_{j\geq1}$ is orthonormal in $L^2(\nu)$. Let $\widehat g$
be any estimated square-integrable payoff rule and write
\[
\widehat a_j
:=
\langle\widehat g,\psi_j\rangle_{L^2(\nu)}.
\]
Decode the signs by
\[
\widehat\theta_{j-M}
:=
\operatorname{sign}(\widehat a_1\widehat a_j),
\qquad M<j\leq2M,
\]
with any fixed convention at zero. This decoder is invariant to any nonzero
global rescaling of $\widehat g$.

Define the scale-free directional error
\[
D_\theta(\widehat g)
:=
\inf_{c\in\mathbb R}
\|g_\theta-c\widehat g\|_{L^2(\nu)}^2.
\]
For any $c\in\mathbb R$, if $c\widehat a_1\leq0$, the squared error in the
first coordinate is at least $a^2\mu_1$. If $c\widehat a_1>0$, every
incorrectly decoded sign contributes at least $h^2\mu_j/M$ to the squared
error. Since $h\leq a$ and $H(\widehat\theta,\theta)\leq M$, both cases
imply
\[
D_\theta(\widehat g)
\geq
\frac{h^2c_\mu}{2^bM^{b+1}}
H(\widehat\theta,\theta).
\]
For every neighboring pair, the sum of the two sign-error probabilities is at
least $1-\mathrm{TV}\geq3/4$. Averaging over the hypercube and summing over
its $M$ coordinates therefore gives
\[
2^{-M}\sum_\theta
\E_\theta[H(\widehat\theta,\theta)]
\geq
\frac{3M}{8}.
\]
Consequently,
\[
2^{-M}\sum_\theta
\E_\theta[D_\theta(\widehat g)]
\geq
\frac{3h^2c_\mu}{2^{b+3}}M^{-b}.
\]

\medskip
\noindent\textit{Step 5: directional error is a Sharpe-ratio gap.}
Write
\[
m:=\langle g_\theta,\widehat g\rangle_{L^2(\nu)},
\qquad
v:=\|\widehat g\|_{L^2(\nu)}^2.
\]
If $v>0$, the payoff generated by $\widehat g$ satisfies
\[
\E_\theta[R(\widehat g)]=m,
\qquad
\E_\theta[R(\widehat g)^2]=v,
\]
and hence
\[
\SR_\theta(\widehat g)^2
=
\frac{m^2}{v-m^2}.
\]
Moreover,
\[
D_\theta(\widehat g)
=
A_M-\frac{m^2}{v}
=
\frac{1}{1+\SR_\theta(\widehat g)^2}
-
\frac{1}{1+(\SR_\theta^\star)^2}.
\]
The same identity holds for $v=0$ under the convention
$\SR_\theta(\widehat g)=0$. Cauchy--Schwarz implies
$|\SR_\theta(\widehat g)|\leq\SR_\theta^\star$.

If $\SR_\theta(\widehat g)\geq0$, then
\[
\SR_\theta^\star-\SR_\theta(\widehat g)
\geq
\frac{D_\theta(\widehat g)}{2\overline S}.
\]
If $\SR_\theta(\widehat g)<0$, then the gap is at least $\underline S$;
since $D_\theta(\widehat g)\leq A_M\leq1/4$, it is again bounded below by
a fixed positive multiple of $D_\theta(\widehat g)$. Therefore there is a
constant $c_S>0$, independent of $M,T,\theta$, such that
\[
\SR_\theta^\star-\SR_\theta(\widehat g)
\geq
c_S D_\theta(\widehat g).
\]
Combining this inequality with Step 4 yields
\[
2^{-M}\sum_\theta
\E_\theta\!\left[
\SR_\theta^\star-\SR_\theta(\widehat g)
\right]
\geq
c_1M^{-b}
\]
for a constant $c_1>0$. Since
$M=\lceil T^{1/(b+1)}\rceil\leq2T^{1/(b+1)}$ for all sufficiently large
$T$, we obtain
\[
\sup_\theta
\E_\theta\!\left[
\SR_\theta^\star-\SR_\theta(\widehat g)
\right]
\geq
c_2T^{-b/(b+1)}
\]
for some $c_2>0$.

The argument allowed $\widehat g$ to range over all square-integrable
estimators in the more informative augmented experiment, a larger class than
managed portfolio rules based on the observed financial rows. Taking the
infimum over all measurable portfolio-learning procedures proves
\[
\inf_{\widetilde W_T}
\sup_{\varepsilon\in\mathcal E}
\mathbb E_\varepsilon
\left[
\SR^\star_{\varepsilon,K}
-
\SR_\varepsilon\!\left(R^p_{t+1}(\widetilde W_T)\right)
\right]
\geq
c\,T^{-b/(b+1)}
\]
for a constant $c>0$ and all sufficiently large $T$.
\end{proof}
